\chapter{Introduction}

Logical reasoning plays a central role in the way humans justify actions,
choices, and conclusions, particularly when decisions are not primarily driven
by emotions or other non-rational factors.

For example, if ``All men are mortal'' and ``Socrates is a man'', then we can
conclude that ``Socrates is mortal''. This form of reasoning, known as an
Aristotelian syllogism, illustrates how conclusions can be derived
mechanically from a collection of premises. A similar principle underlies
mathematics, where the validity of a statement is established through a proof
constructed from axioms and inference rules.

The search for a formal foundation for mathematics dates back to the work of
Frege, whose \emph{Begriffsschrift} (1879) attempted to express mathematical
reasoning within a single logical formalism. Russell's paradox later revealed
that Frege's system was inconsistent, motivating the development of new
foundational frameworks. Today, several approaches coexist, including
axiomatic set theory, intuitionistic logic, and various forms of type theory.

In parallel, Church introduced the $\lambda$-calculus as a mathematical model
of computation~\cite{church}. Originally untyped, it was later enriched with
types to prevent non-terminating computations and eventually extended with
polymorphism, inductive definitions, and dependent types. These developments
culminated in the \emph{Calculus of Inductive Constructions} (CIC~\cite{cic}), 
which combines a dependently typed
$\lambda$-calculus with inductive types. CIC forms the logical foundation of
the \rocq proof assistant (formerly known as \coq).

A fundamental connection between logic and computation was made by Curry and
Howard~\cite{Howard-correspondance}, who observed that logical propositions
and types share the same structure. Under this
\emph{Curry--Howard correspondence}, propositions are interpreted as types,
proofs as programs, and proof checking becomes ordinary type checking.

A simple example is the proposition \textit{A implies A}, that we can write,
$ A \rightarrow A. $
Under the Curry--Howard correspondence, this proposition is interpreted as the
type of functions taking an argument of type \(A\) and returning a result of
the same type. The identity function
$ \lambda (x:A).x $
is therefore simultaneously a program and a proof of the proposition
\(A \rightarrow A\). More generally, constructing a proof amounts to writing a
well-typed program, while checking the correctness of the proof reduces to
verifying that the corresponding program typechecks.

An essential feature of CIC is the support for \emph{inductive types},
which provide a principled mechanism for defining recursive data structures.
For example, the natural numbers are defined as
\cI{nat := O | S : nat ~$\rightarrow$~ nat },
while polymorphic lists are defined as
\cI{list A := nil | cons : A ~$\rightarrow$~ list A ~$\rightarrow$~ list A.}

Every natural number is therefore obtained by repeatedly applying the
constructor \cI{S} to \cI{O}, whereas every list is either empty or consists of
a head element followed by another list. Such definitions automatically give
rise to recursion and induction principles, making inductive types a central
mechanism for programming and proving in \rocq.

Using software to construct proofs has become an increasingly viable
alternative to pen-and-paper mathematics. Rather than trusting the human
author, one only needs to trust the small kernel of the proof assistant,
whose role is to verify that every proof term is well typed. Once a proof has
been accepted by the kernel, the desired theorem is guaranteed to follow from
the underlying logical foundations.

For example, addition over natural numbers can be defined recursively, after
which one may prove that zero is a right neutral element:
$ \forall n : \cI{nat},\; n + 0 = n. $
Although the proof proceeds by induction over the structure of natural
numbers, the proof assistant mechanically checks every inference step,
eliminating the possibility of mistakes in the logical argument.

The importance of machine-checked proofs is particularly evident in
safety-critical domains such as aerospace engineering, medicine, and
cryptography, where software errors may have catastrophic consequences. A
well-known example is the failure of the Ariane~5 maiden flight in 1996,
caused by a software defect.

Beyond checking proofs, modern proof assistants also assist users in
constructing them. Although every proof ultimately has a fully explicit
representation, requiring users to write this representation directly would
make proof development prohibitively verbose. Instead, systems such as
\rocq allow users to omit information that can be reconstructed
automatically.

These mechanisms are collectively referred to as
\emph{elaboration}~\cite[Chap.~6]{elaboration}. The elaborator bridges the gap
between the concise notation written by users and the fully explicit proof
terms manipulated by the kernel. Besides expanding notations, elaboration
infers omitted arguments, reconstructs implicit types, inserts coercions,
resolves overloaded symbols, and generates proof obligations left implicit by
the user.

For instance, the list
\cI{[1;2;3]}
is internally represented by the considerably more explicit term
\cI{cons nat 1 (cons nat 2 (cons nat 3 (nil nat)))}.

A common source of ambiguity arises from symbol overloading. In mathematics,
the notation \(+\) may denote integer addition, vector addition, or matrix
addition, with the intended meaning inferred from the surrounding context.
Similarly, in \rocq, a notation may admit several implementations. The
elaborator must therefore infer the appropriate interpretation by solving a
collection of constraints generated during type checking.

As observed in~\cite[Chap.~6]{elaboration}, elaboration algorithms are often
specified using inference rules strongly reminiscent of logic programming.
Indeed, resolving an elaboration problem, such as finding a missing
argument, an implicit coercion, or a type-class instance, amounts to
searching for a term that satisfies a given goal, drawn from a database of
candidate rules: user-declared instances, coercions, or notations. This
search may require exploring several candidates and backtracking when a
choice turns out to be unsuitable, which is precisely the computational
model implemented by logic-programming interpreters through unification and
backtracking over a database of clauses. This structural correspondence
suggests that a logic-programming language is not merely a convenient
implementation vehicle for elaboration, but a reasonably natural model of
the search process elaboration already performs. In this thesis, we
investigate an alternative implementation of \rocq's type-class resolution
mechanism based on the logic programming language \elpi. This approach
makes it possible to reuse mature logic-programming engines and techniques
instead of reimplementing a dedicated inference procedure inside the proof
assistant. At the same time, it raises several questions. How naturally can
\rocq's inference problems be translated into a logic-programming setting?
What are the performance implications of such a translation? More
generally, what are the benefits and limitations of this design choice?

Beyond type-class resolution, the intrinsic nature of logic programming raises
another important concern, especially for users unfamiliar with logic
programming. Controlling the backtracking performed by the interpreter is
fundamental to managing performance. This motivates a second line of work
developed in this thesis: a determinacy analysis for \elpi
programs~(\cref{chap:det-elpi}). Determinacy ensures that selected predicates
produce at most one solution and do not introduce unexpected choice points for
any call. This property helps control the behavior of logic programs. In our
setting, type classes correspond to predicates and instances to rules. If a
class is declared deterministic, the analysis ensures that the program generated
by the instances of this class satisfies the property that any call to the class
produces at most one solution.

We design and implement an analysis based on inspecting clauses and tracking
argument modes. The algorithm supports most features of \elpi, including
higher-order predicates, dynamic clauses, and the cut operator. Handling cut is
important, as it can enforce determinism while interacting in subtle ways with
logical semantics.

To increase confidence in the static analysis, we formalize the expected
properties of the checker in the \rocq proof assistant by encoding the \elpi
interpreter~(\cref{chap:mechanization}). To simplify this development, we restrict
the model to a first-order fragment of the language, where terms do not contain
higher-order variables. In this setting, the absence of higher-order variables
simplifies the formalization and allows us to produce a fully mechanized proof.

Although this model does not capture all higher-order features of \elpi, it
encompasses the aspects relevant to our analysis and provides a first foundation
for establishing machine-checked correctness guarantees for the \elpi language.

The remainder of this introduction presents the contributions of this thesis
towards answering these questions.


\section{Contributions and structure of the thesis}

This thesis is organized as follows. We first present the
context of the work (\cref{chap:context}). In particular, we describe the
type-class resolution mechanism implemented in \rocq, together with its main
options. Then, we introduce the \elpi language and give a formal description of
its syntax and semantics. Finally, we present the \rocqelpi plugin, which enables
the interaction between \elpi and \rocq.

This thesis contains five main contributions, organized as follows:

\begin{description}
  \item [\Cref{sec:elpi-syntax,sec:elpi-semantics}] We present a formalization
            of the \elpi language, covering both its syntax and its semantics.
            This formalization first appeared in the appendix
            of~\cite{fissore2026}. We include it in \cref{chap:context}, the
            chapter describing the context of this work, because this
            placement fits the overall structure of the thesis better.

  \item[\Cref{chap:tc}] We study the implementation of a type-class solver
                        written in \elpi for \rocq~\cite{fissore2023}. We present the core part of
                        the compiler that translates \rocq classes and
                        instances into an \elpi database. This database is then
                        used by the interpreter to search for inhabitants of
                        type classes. We also compare our solver with \rocq's
                        native solver, focusing on performance, modes, and
                        unification issues.

  \item[\Cref{chap:ho-for-free}] We analyze unification problems arising from
                the higher-order abstract syntax used to represent \rocq terms
                in \elpi. In \rocq, terms that are equal up to
                $\beta\eta$-conversion can unify, while in \elpi such
                unification may fail, even if its algorithm supports these
                conversions. Based on our work~\cite{fissore2024}, we explain how
                to address this issue by delaying problematic unifications using
                a system of links.

  \item[\Cref{chap:det-elpi}] We present a static determinacy checker
               implemented in \elpi, available since version 3. This tool
               verifies that predicates declared as deterministic by the user
               effectively behave deterministically, meaning that they never
               create additional choice points. This result is based on our
               work~\cite{fissore2026}.

  \item[\Cref{chap:mechanization}] We mechanize the determinacy checker in the
               \rocq proof assistant. The formalization considers a restricted
               fragment of \elpi, namely its first-order part, where variables
               only represent data. Within this setting, we prove the main
               determinacy property: in a program that satisfies the checker,
               any call to a deterministic predicate produces at most one
               solution.
\end{description}

Finally, \cref{app:builtins} provides a list of the \elpi built-ins and auxiliary
predicates used in the code snippets throughout this thesis.